\chapter{System identification, recording and flight analysis}\label{ch:ident}
\sloppy

The EKF, the position hold of the driver, the zone guard, the missions and the fake drone of the simulator all
rely on the same few numbers: how fast the Tello really goes for a given velocity command, how quickly it gets
there, and how late. This chapter explains where those numbers come from and how the project checks, after
every flight, that the whole chain still behaves. Six programs are involved:

\begin{center}
\small
\begin{tabular}{@{}lll@{}}
\toprule
Program & Role & Main output \\
\midrule
\file{tello\_driver/step\_test.py} & identification / validation flight & \file{step\_<time>/*.csv} \\
\file{tello\_driver/step\_fit.py} & fits $\tau \dot v = K u(t-L) - v$ per axis & \file{model.yaml}, plots \\
\file{tello\_driver/session\_recorder.py} & records every flight (ROS node) & \file{sessions/<time>/} \\
\file{tools/flight\_report.py} & after-flight metrics and verdicts & \file{report.txt}, \file{report.png} \\
\file{tools/calibrate\_camera.py} & checkerboard calibration & \file{tello\_camera.yaml} \\
\file{tools/make\_print\_sheets.py} & printable markers and checkerboard & \file{docs/print/*.pdf} \\
\bottomrule
\end{tabular}
\end{center}

The identified constants are copied by hand into \file{tello\_driver/model.py}, ``the ONE place where the
model lives''. Section~\ref{sec:ident-use} shows every place that reads them.

% =====================================================================================================
\section{The model being identified}\label{sec:ident-model}

The Tello hides its motors and its own flight controller: the only input we have is the \code{rc} velocity
command (sent at 20\,Hz by the driver), and the only accurate output is the ArUco pose. So the project identifies
the \emph{closed loop} ``Tello + its own controller'', from the command $u$ (m/s or rad/s) to the velocity $v$
actually reached, per axis (forward $x$, left $y$, up $z$, yaw), as a first order system with a pure delay:
$$
\tau \frac{dv}{dt}(t) = K\,u(t-L) - v(t), \qquad p(t) = p(0) + \int_0^t v(s)\,ds .
$$
\begin{itemize}[nosep]
\item $K$ (dimensionless): steady-state ratio real speed / commanded speed;
\item $\tau$ (s): time constant of the velocity response;
\item $L$ (s): total delay from the command to the \emph{marker pose}, i.e.\ $L = L_u + L_v$ with $L_u$ the
command side (Wi-Fi, flight controller) and $L_v$ the video side (image transport and detection).
\end{itemize}

\paragraph{Step response.} For a unit step of $u$ at $t=0$ and $v(0)=0$, the velocity is
$v(t) = K(1-e^{-(t-L)/\tau})$ for $t>L$ and the displacement is
$$
p(t) - p(0) = K\Big[(t-L) - \tau\big(1 - e^{-(t-L)/\tau}\big)\Big] \xrightarrow[t \gg L+\tau]{} K\,\big(t - (L+\tau)\big).
$$
After the transient the position is a ramp of slope $K$ that is \emph{late by exactly $L+\tau$}.

\begin{intuition}
A delay $L$ and a lag $\tau$ look alike in a position trace: both make the drone ``arrive late''. Only the
curvature of the first few tenths of a second distinguishes them, and that part is buried in the
$\sim$1--2\,cm noise of the ArUco pose. Hence the fit determines $L+\tau$ well and $L$, $\tau$ separately
poorly: in the $(\tau, L)$ plane the cost has a long valley along the line $L+\tau = \text{const}$.
This is why \file{step\_fit.py} draws \file{cost\_map.png} and reports the spread of $L+\tau$, and why
\file{model.py} chose $\tau_{xy}=0.25$\,s ``to keep the measured $L+\tau$''.
\end{intuition}

\paragraph{The excitation: doublets.} A doublet of amplitude $A$ and half-duration $T$ is $u=+A$ during $T$,
then $u=-A$ during $T$. Ideally the drone goes $\approx KAT$ away and comes back, so it stays in a small box,
and the two edges give two transients per test. Several $(A,T)$ pairs are used because, as the source notes
(Nov 13), two identical doublets ``could not separate the delay $L$ from the time constant $\tau$''.

\paragraph{Three design decisions} recur in the code, each motivated by a real failure:
\begin{enumerate}[nosep]
\item \textbf{Input = the command really sent} (\topic{/tello/rc\_sent}, after the driver's guards), not the
command asked. On Nov 13 the height-band guard cut an up pulse from 0.30 to 0.08\,m/s and the fit on the
requested command found $K_z = 0.43$.
\item \textbf{Output = raw ArUco position} (\topic{/aruco/pose}), not the EKF (which itself uses the model
being identified: circular) and not a velocity (differentiating a noisy pose amplifies the noise).
\item \textbf{Per-doublet nuisance terms} $c_j + d_j(t-t_j)$: an offset and a linear drift absorb where the
doublet starts and the Tello's slow drift ($\sim$10\,cm/s on its own, Nov 10).
\end{enumerate}

% =====================================================================================================
\section{The identification flight: \texorpdfstring{\texttt{step\_test.py}}{step\_test.py}}\label{sec:steptest}

Run with \code{ros2 run tello\_driver tello\_step\_test} while \code{real.launch.py} runs in another terminal
(the driver, localizer and EKF must be up). The parameter \code{schedule:=valid} flies a shorter
\emph{validation} flight with different doublets. The class \code{StepTest} is also the base class of the
missions (\code{MarkerTour} in \file{marker\_tour.py} inherits it), which explains a few hooks
(\code{fence}, \code{MOVE\_TO\_ZONE\_CENTRE}, \code{POSE\_TOPIC}, \code{PREFIX}) that the identification itself
does not use.

\subsection{Module level}

\code{RATE = 20.0} Hz is the control loop rate (the driver's rate too). \code{SCHEDULES} holds the
$(A, T)$ pairs, in m/s (rad/s for yaw) and seconds:

\begin{center}
\small
\begin{tabular}{@{}lll@{}}
\toprule
 & \code{ident} & \code{valid} \\
\midrule
x, y & (0.15,1.8) (0.25,1.2) (0.35,0.8) & (0.20,1.0) (0.30,1.5) \\
 & (0.25,1.8) (0.35,1.2) (0.15,0.8) & (0.30,0.6) (0.20,1.5) \\
z & (0.30,0.4) (0.25,0.5) (0.30,0.3) (0.25,0.4) & (0.28,0.4) (0.30,0.35) \\
yaw & (0.3,1.2) (0.5,1.0) (0.7,0.6) (0.5,0.6) & (0.4,0.8) (0.6,1.0) \\
\bottomrule
\end{tabular}
\end{center}

With the default \code{axes = x,y,z,yaw}, the identification flight has $6+6+4+4 = 20$ doublets. Each lasts
$2 + 2T + 3$\,s (hover, pulses, hover), i.e.\ $\sum (5 + 2T) = 45.2 + 45.2 + 23.2 + 26.8 = 140.4$\,s, plus a
re-centring of up to 7\,s before each: the ``$\sim$3.5\,min'' of the docstring. The vertical doublets are short
but strong: $AT \le 0.125$\,m so that the band guard (band $1.70$--$2.00$\,m) does not cut them, and
$A \ge 0.25$\,m/s because the Tello ``did not react at all to 8--12 rc points up/down''.

\func{Abort}{(Exception)}
Raised for every reason to stop the test (key pressed, low battery, out of the box, \dots). Caught in
\code{run()}, which then lands.

\func{Lost}{(Exception)}
Raised when the markers are lost for more than 1\,s \emph{during a doublet}: the doublet is abandoned, its
segment number is marked as failed, and it is repeated (up to 3 attempts).

\func{wrap}{(a) -> float}
Wraps an angle into $(-\pi, \pi]$ with $\operatorname{atan2}(\sin a, \cos a)$.

\func{driver\_zone\_file}{(node, timeout=5.0) -> str}
Asks the running driver which flight-zone file it enforces, so that the test and the driver agree. It creates a
client of the standard parameter service of the driver node, \topic{/tello\_driver/get\_parameters}, of type
\code{rcl\_interfaces/srv/GetParameters}. It waits up to 5\,s for it, requests the single parameter
\code{flight\_zone\_file}, spins until the future completes (5\,s again), and returns
\code{values[0].string\_value}. Any failure (no driver, timeout, exception) returns \code{''}; the client is
destroyed in a \code{finally}.

\subsection{Construction}

\func{StepTest.\_\_init\_\_}{(self)}
Declares the parameters and opens the log files. In order:
\begin{itemize}[nosep]
\item \code{schedule} (\code{ident}/\code{valid}, anything else exits), \code{axes} (comma-separated list),
\code{box = 0.8} m (half-size of the horizontal safety box around the start point), \code{z\_min = 0.4},
\code{z\_max = 2.4} m, \code{lost\_timeout = 15} s (blind for that long: land).
\item \textbf{Flight zone.} \code{band} defaults to \code{True}. The zone file is, by priority: the parameter
\code{flight\_zone\_file}; else the driver's, via \code{driver\_zone\_file()}; else
\file{\$TELLO\_REPO/config/flight\_zone\_band.yaml} (or \file{flight\_zone.yaml} when \code{band:=false}). It is
loaded with \code{zone.load\_zone()} (rectangle, margin, free-space grid). In band mode a missing zone is fatal;
the band $(z_{\min}, z_{\max}) = (1.70, 2.00)$ is read from it and the abort limits become the band
$\pm 0.3$\,m, i.e.\ $[1.40, 2.30]$\,m.
\item \code{min\_battery}: 50\,\% for the identification (\code{PREFIX == "step"}), 30\,\% for subclasses
(``at 32\,\% the Tello only climbed to 1.28\,m'').
\item Output folder \file{\textasciitilde/tello\_ws/sessions/step\_YYYYmmdd\_HHMMSS} (\code{stepval\_\dots}
for a validation flight).
\item Publisher on \topic{/tello/cmd\_vel\_safe}, \code{Trigger} clients for \topic{/tello\_driver/takeoff}
and \topic{/tello\_driver/land}, subscriptions to the pose (\topic{/aruco/pose} by default, parameter \code{pose\_topic}),
\topic{/tello/telemetry}, \topic{/tello/camera\_info} (only used as a ``video is alive'' heartbeat) and
\topic{/tello/rc\_sent} (written straight to \file{sent.csv} by a lambda).
\end{itemize}

The four CSV files of the step folder:

\begin{center}
\begin{tabular}{@{}lll@{}}
\toprule
File & Columns & Written by \\
\midrule
\file{poses.csv} & \code{t, x, y, z, yaw, t\_image} (epoch s, m, rad) & \code{on\_pose} \\
\file{commands.csv} & \code{t, segment, axis, u, cx, cy, cz, cw} & \code{send} \\
\file{sent.csv} & \code{t, cx, cy, cz, cw} (what the driver really sent) & \topic{/tello/rc\_sent} lambda \\
\file{telemetry.csv} & every key of the telemetry JSON, plus \code{t} & \code{on\_tel} \\
\file{failed\_segments.txt} & segment numbers abandoned (markers lost) & \code{close} \\
\bottomrule
\end{tabular}
\end{center}

All times \code{t} are \code{time.time()} at \emph{reception}; \code{t\_image} is the header stamp of the pose
(time the image was taken). \code{segment} is $-1$ outside the doublets.

\subsection{Inputs}

\func{on\_pose}{(self, m)}
Rejects the pose when $\sqrt{\max(\Sigma_{xx}, \Sigma_{yy})} > 0.15$\,m (\code{max\_pose\_sigma}); this matters
for the missions that subscribe to the EKF, whose covariance grows without markers. Extracts the yaw from the
quaternion, $\psi = \operatorname{atan2}(2(q_wq_z + q_xq_y),\ 1 - 2(q_y^2+q_z^2))$, stores
\code{self.pose = (t\_received, x, y, z, yaw)}, updates the offset ArUco yaw minus IMU yaw
(\code{yaw\_offset}, used when the markers disappear) and writes one row of \file{poses.csv}.

\func{on\_tel}{(self, m)}
Parses the telemetry JSON (ignores malformed messages), keeps it in \code{self.tel} with its reception time, and
appends it to \file{telemetry.csv}; the CSV header is the key set of the first message.

\func{on\_video}{(self, m)}
Records the time of the last \code{camera\_info} (one per frame): if it is older than 1\,s the Wi-Fi video is
considered lost.

\func{imu\_yaw}{(self) -> float or None}
The Tello reports \code{yaw} in clockwise-positive degrees; the map convention is counter-clockwise radians, so
it returns $-\text{yaw}\cdot\pi/180$ (checked on 4 sessions: slope $-0.93$ to $-1.08$). \code{None} if the
telemetry is older than 0.5\,s.

\func{yaw\_estimate}{(self) -> float or None}
The ArUco yaw when the pose is younger than 0.5\,s, else $\operatorname{wrap}(\psi_{\text{IMU}} + \text{offset})$.

\func{in\_air}{(self) -> bool}
\code{flying} in the telemetry, or a height \code{h} $\ge$ 20\,cm.

\func{pose\_age}{(self) -> float}
Seconds since the last accepted pose ($\infty$ if none).

\subsection{Helpers and the built-in safety}

\func{say}{(self, text)}
Prints with carriage returns around (the terminal is in cbreak mode).

\func{spin}{(self, seconds)}
\code{spin\_once} with a 20\,ms timeout until \code{seconds} have passed: processes callbacks while waiting.

\func{call}{(self, name, timeout=5.0) -> Trigger.Response}
Calls \topic{/tello\_driver/takeoff} or \topic{/tello\_driver/land}; raises \code{Abort} if the service does not appear within
2\,s.

\func{check\_keys}{(self)}
Non-blocking \code{read\_key(0)} (from \file{teleop.py}); any key returns a value and raises
\code{Abort("key pressed")}.

\func{safety}{(self)}
Called at every control tick, independently of the driver's own safety chain (defence in depth). It raises
\code{Abort} when:
\begin{enumerate}[nosep]
\item a key is pressed;
\item the battery is below 20\,\%;
\item with a fresh pose ($<0.5$\,s) and a start point: the free-space clearance (from the grid) is below
\code{margin} $- 0.2$\,m (with the band zone's margin 0.25\,m: closer than 5\,cm to a scanned obstacle); or,
without grid, the drone is more than 0.2\,m outside the zone rectangle;
\item the drone leaves the \code{fence} rectangle $\pm$ \code{fence\_tol} (missions only), or else the box
$|x - x_0| > 0.8$ or $|y - y_0| > 0.8$\,m around the start point $p_0$;
\item the height leaves $[z_{\min}, z_{\max}]$ \emph{after having once been inside} (\code{z\_reached}): a weak
climb must not abort the test on the way up;
\item no pose for more than \code{lost\_timeout} (15\,s).
\end{enumerate}

\func{send}{(self, cx=0, cy=0, cz=0, cw=0, segment=-1, axis="", u=0, blind=False)}
Publishes a \code{Twist} on \topic{/tello/cmd\_vel\_safe} (body frame: \code{linear.x} forward, \code{linear.y}
left, \code{linear.z} up, \code{angular.z} yaw rate) and logs it in \file{commands.csv}. If the pose is older
than 0.7\,s and \code{blind} is false, everything (including the logged \code{u}) is zeroed: the drone hovers
rather than continuing a test blind.

\func{search\_heading}{(self, base\_yaw)}
Markers out of view: yaw-only search (the driver blocks horizontal motion without a pose anyway). For a loss of
$\ell$ seconds the target heading is
$\psi^\star = \psi_0 + 40^\circ \sin\big(2\pi(\ell-3)/6\big)$ once $\ell>3$\,s (a $\pm40^\circ$ sweep of period
6\,s), and the command is $\omega = \operatorname{clip}(0.9\,\operatorname{wrap}(\psi^\star - \hat\psi), \pm 0.5)$\,rad/s
sent with \code{blind=True}, $\hat\psi$ from \code{yaw\_estimate()}.

\func{search\_until\_found}{(self, base\_yaw, timeout=None) -> bool}
Loops \code{safety()} + \code{search\_heading()} at 20\,Hz until a pose younger than 0.3\,s arrives (\code{True})
or the optional timeout expires (\code{False}). The overall limit is the 15\,s \code{lost\_timeout} of
\code{safety()}.

\func{run}{(self, seconds, fn, restart\_on\_loss=False, search\_on\_loss=True)}
The control loop: for \code{seconds}, at 20\,Hz on a fixed schedule (\code{nxt += 1/RATE}, spinning in the
remaining time), call \code{safety()} then \code{fn()}; \code{fn()} returning \code{False} ends the loop early.
If the pose is older than 1\,s (and a start point exists): during a doublet (\code{restart\_on\_loss}) raise
\code{Lost}; otherwise search on the spot and resume.

\subsection{Flight phases}

\func{recentre}{(self, timeout=6.0, target=None)}
Proportional return to $p_0$ (or \code{target}) in position and heading. With the world errors
$e_x, e_y, e_z, e_\psi$ and the current yaw $\psi$, the horizontal error is rotated into the body frame,
$$
\begin{bmatrix} b_x \\ b_y \end{bmatrix} =
\begin{bmatrix} \cos\psi & \sin\psi \\ -\sin\psi & \cos\psi \end{bmatrix}
\begin{bmatrix} e_x \\ e_y \end{bmatrix},
$$
and the command is $(\operatorname{clip}(0.6b_x, 0.25), \operatorname{clip}(0.6b_y, 0.25),
\operatorname{clip}(0.6e_z, 0.25), \operatorname{clip}(0.8e_\psi, 0.4))$. It stops when
$\lVert e_{xy}\rVert < 8$\,cm, $|e_z|<8$\,cm and $|e_\psi|<5^\circ$, then sends zero for 1\,s.

\func{return\_home}{(self)}
After the tests: flies back above the takeoff pad of \file{config/home.yaml} at the current height (timeout
15\,s), so that the landing ends on it. Falls back to \code{recentre()} if the pad is unknown, the pose is
stale, or the pad's clearance in the free-space grid is below the margin.

\func{hold\_cmd}{(self) -> dict}
The ``other axes hold the position'' command used during the doublets: same structure as \code{recentre}
toward $p_0$ with gains 0.6 (position) and 0.8 (yaw) but tighter limits (0.2\,m/s, 0.3\,rad/s), and zero
without a fresh pose. Needed because the Tello drifts by itself (37\,cm in 4\,s of ``hover'' on Nov 10).

\func{doublet}{(self, segment, axis, sign, amp, dur)}
The core of the experiment. Four phases run through \code{run(\dots, restart\_on\_loss=True)}:
\begin{lstlisting}
for u, d in ((None, 2.0), (sign * amp, dur), (-sign * amp, dur), (None, 3.0)):
    self.run(d, lambda u=u: cmd(u), restart_on_loss=True)
\end{lstlisting}
At each tick \code{cmd(u)} takes \code{hold\_cmd()} for all axes and overrides the tested one with $u$ when it is
not \code{None}: during the pulses the tested axis is \emph{open loop}, during the two hovers it is held like the
others. The logged \code{u} is the value of the tested axis in both cases, with the \code{segment} number and the
axis. (Note the default argument \code{u=u}: it freezes the loop variable inside the lambda.)

\func{wait\_ready}{(self)}
Waits up to 10\,s for telemetry, then checks: no other publisher on \topic{/tello/cmd\_vel\_safe} (via
\code{get\_publishers\_info\_by\_topic}: a running teleop would mix its commands into the test), video alive
(a \code{camera\_info} less than 1\,s old), battery $\ge$ \code{min\_battery}. Each failure is an \code{Abort}.

\func{main}{(self)}
\code{prepare()} then \code{run\_tests()}.

\func{prepare}{(self)}
\begin{enumerate}[nosep]
\item \code{wait\_ready()}. If on the ground: wait for ENTER (any other key cancels), call takeoff, then wait up
to 25\,s until the telemetry says flying and not \code{busy} (the driver takes off, then climbs to the band
itself).
\item \textbf{Start point}: collect up to 15 poses younger than 0.2\,s (within 30\,s). If no marker is seen after
2\,s, turn on the spot at 0.5\,rad/s (\code{blind}) until some appear. Fewer than 5 samples: abort. $p_0$ is the
per-coordinate \emph{median} of the samples (robust to a wrong pose).
\item \textbf{Move to the centre} (\code{MOVE\_TO\_ZONE\_CENTRE}): the target is the middle of the zone rectangle,
$((x_{\min}+x_{\max})/2, (y_{\min}+y_{\max})/2) = (0.04, 1.56)$\,m for the band zone, at the middle of the band
(1.85\,m), via \code{recentre(timeout=12)}. Then $p_0$ becomes the reached pose (with $z_0 = 1.85$).
\item Height check: $p_0$ inside the abort limits (band mode), or 0.3\,m inside them otherwise.
\end{enumerate}

\func{run\_tests}{(self)}
Builds the plan \code{[(axis, A, T)]} from the schedule, axis by axis. The sign alternates per axis
(+ for the 1st, 3rd, \dots\ doublet of an axis, $-$ otherwise) so that both directions are excited. For each
doublet, up to 3 attempts: \code{recentre()}, then \code{doublet()}; on \code{Lost} the segment number goes to
\code{failed\_segments}, a new number is used for the retry, and the drone searches for the markers. Segment
numbers are therefore unique and \file{step\_fit.py} can simply skip the failed ones. Ends with
\code{return\_home()}.

\func{close}{(self)}
Closes the four CSV files and writes \file{failed\_segments.txt} (space-separated numbers).

\subsection{Entry points}

\func{main}{() -> None}
\code{run(StepTest, fit\_after=True)}.

\func{run}{(cls, fit\_after=False)}
Shared by the step test and the missions. Forces local-only DDS discovery
(\code{ROS\_AUTOMATIC\_DISCOVERY\_RANGE=LOCALHOST}, \code{local\_ros\_env()}), initialises rclpy with
\code{SignalHandlerOptions.NO} so that Ctrl+C raises \code{KeyboardInterrupt} in Python instead of shutting
rclpy down (the node must still be able to call \emph{land}), prints paragraphs 3--4 of the module docstring,
puts the terminal in cbreak mode, and runs \code{node.main()}. Whatever happens (\code{Abort}, Ctrl+C, normal end)
the \code{finally} publishes a zero \code{Twist}, calls land if \code{in\_air()}, restores the terminal and closes
the files. Only after a \emph{successful} run with \code{fit\_after}: an identification flight is fitted at once
(\code{step\_fit.fit}), a validation flight is validated against the latest earlier identification
(\code{step\_fit.validate}).

\begin{keypoint}
The test has three safety layers: (1) the driver's chain (\topic{/tello/cmd\_vel\_safe} is filtered: hold, ceiling,
band, zone, wall), whose output is exactly what \file{sent.csv} records; (2) the test's own \code{safety()} with
its box, height limits and blind timeout; (3) any key or Ctrl+C lands. The driver's guards may modify a pulse,
which is why the fit uses \file{sent.csv}.
\end{keypoint}

% =====================================================================================================
\section{Fitting the model: \texorpdfstring{\texttt{step\_fit.py}}{step\_fit.py}}\label{sec:stepfit}

\code{ros2 run tello\_driver tello\_step\_fit [folder]} fits the latest (or the given) \file{step\_*} folder;
\code{--validate <ident> <valid>} simulates an identified model on another flight. Module constants: the axes
\code{x, y, z, yaw} and their column in \file{sent.csv} (\code{cx, cy, cz, cw}), the simulation step
$\Delta = $ \code{DT} $= 0.01$\,s, the search grids
\begin{center}
\begin{tabular}{@{}lll@{}}
\toprule
grid & $\tau$ (s) & $L$ (s) \\
\midrule
fine & $0.05:0.025:1.5$ (59 values) & $0:0.01:0.9$ (91 values) \\
coarse (leave-one-out) & $0.05:0.05:1.5$ (30 values) & $0:0.02:0.9$ (46 values) \\
\bottomrule
\end{tabular}
\end{center}
and \code{MODEL\_IN\_USE}, the reference model the validation compares with (see the review).

\subsection{Data}

\func{rows}{(path) -> list[dict]}
\code{csv.DictReader} rows, or \code{[]} if the file does not exist.

\func{load}{(folder) -> (segs, poses, has\_sent)}
Reads \file{poses.csv} into an array $[t, x, y, z, \psi]$; reads \file{commands.csv} and groups the rows by
\code{segment}, skipping $-1$ and the failed segments, into \code{segs[sid] = \{axis, t, u\}}. If
\file{sent.csv} exists, the asked \code{u} is kept as \code{u\_asked} and replaced by the \emph{really sent}
command on the segment's axis, by zero-order hold: for each command time $t_k$, the last \file{sent.csv} row with
$t \le t_k$ (\code{searchsorted(\dots, side="right") - 1}). \code{has\_sent} tells whether this was possible (older
flights had no \file{sent.csv}).

\func{measured}{(seg, poses) -> (t, y) or (None, None)}
The output signal of one doublet: the pose samples between the first and last command of the segment (at least
10, else \code{None}). The reference heading $\psi_0$ is the median yaw over the first second (the hover), and the
position is projected on the \emph{body axis at the start}:
$$
y_x = x\cos\psi_0 + y\sin\psi_0, \qquad y_y = -x\sin\psi_0 + y\cos\psi_0, \qquad y_z = z, \qquad y_\psi = \operatorname{unwrap}(\psi).
$$
Using a fixed axis is valid because the other axes are held (the yaw hold keeps $\psi \approx \psi_0$).

\func{response}{(t\_cmd, u, t\_eval, tau, L) -> array}
Simulates the position for $K = 1$, $g_{\tau,L}(t)$, sampled at the pose times. On a uniform grid
$t_k = t_0 + k\Delta$ starting at the first command:
\begin{enumerate}[nosep]
\item delayed input by zero-order hold: $\tilde u_k = u(t_k - L)$ (last command at or before $t_k-L$; $0$ before
the first command);
\item forward-Euler discretisation of $\tau\dot v = \tilde u - v$ with $a = \Delta/\tau$:
$$
v_k = (1-a)\,v_{k-1} + a\,\tilde u_{k-1}, \qquad v_0 = 0,
$$
implemented as the IIR filter \code{lfilter([0, a], [1, -(1-a)], ug)} (transfer function $a z^{-1}/(1-(1-a)z^{-1})$),
with a pure-Python loop if SciPy is missing;
\item position by rectangle integration $p_0 = 0$, $p_k = \Delta\sum_{i<k} v_i$, then linear interpolation at
\code{t\_eval}.
\end{enumerate}
The scheme is stable for $a<1$, i.e.\ always on the grid ($\tau \ge 0.05 > \Delta$).

\subsection{Least squares and the grid search}

For an axis with doublets $j = 1..n$, samples $(t_{ji}, y_{ji})$ and simulated responses $g_{ji} = g_{\tau,L}(t_{ji})$,
the model is
$$
y_{ji} \approx K\,g_{ji} + c_j + d_j\,(t_{ji} - t_{j1}).
$$
For \emph{fixed} $(\tau, L)$ it is linear in $\theta = (K, c_1, d_1, \dots, c_n, d_n)$: one shared gain, and one
offset and drift per doublet. Stacking all samples, $A\theta \approx b$ with
$$
A = \begin{bmatrix}
\mathbf g_1 & \mathbf 1 & \mathbf t_1 - t_{11} & & & \\
\mathbf g_2 & & & \mathbf 1 & \mathbf t_2 - t_{21} & \\
\vdots & & & & & \ddots
\end{bmatrix},\qquad b = \begin{bmatrix}\mathbf y_1\\ \mathbf y_2\\ \vdots\end{bmatrix},
$$
and the cost is the residual sum of squares
$$
J(\tau, L) = \min_\theta \lVert A(\tau, L)\,\theta - b\rVert^2 .
$$
This is a \emph{separable} least-squares problem: the linear parameters are eliminated exactly and only the two
nonlinear ones are searched, here exhaustively on the grid, $(\hat\tau, \hat L) = \arg\min J$.

\func{solve}{(data, g, K=None) -> (K, sse, model, meas)}
Builds $A$ and $b$ as above and solves with \code{np.linalg.lstsq}. With \code{K} given (validation), the
$K$ column is dropped and $K g$ is subtracted from $b$: only $c_j, d_j$ are fitted; \code{model} and \code{meas}
are then shifted back by $Kg$ so that both cases return the model and the measured positions. Returns the gain,
$J$ (\code{sse}), and the two stacked vectors.

\func{fit\_percent}{(data, model, meas) -> float}
The normalised-RMSE fit used by MATLAB's \code{compare}:
$$
\text{fit} = 100\left(1 - \frac{\lVert \mathbf y - \hat{\mathbf y}\rVert}{\lVert \mathbf y - \bar{\mathbf y}\rVert}\right),
$$
where $\bar{\mathbf y}$ is the mean \emph{of each doublet} (so the reference is ``a constant per doublet'').
100 is perfect, 0 is no better than a constant, negative is worse.

\func{grid\_fit}{(data, taus, Ls, cache=None) -> (K, tau, L, J)}
Double loop over the grid: simulate every doublet, \code{solve}, store $J$ in a map of size
\code{len(taus)} $\times$ \code{len(Ls)} ($59 \times 91 = 5369$ evaluations per axis on the fine grid), keep the
minimum. The optional \code{cache} memoises the simulated responses with the key
\code{(id(seg), i, k, len(taus))}: during leave-one-out the same doublet is simulated at the same grid point in
every subset, so each response is computed once.

\func{fit\_axis}{(segs, poses) -> dict or None}
Needs at least two usable doublets. Fine-grid fit, then the final \code{solve} at the optimum (for the RMS and the
fit \%). \textbf{Uncertainty} by leave-one-out on the coarse grid: refit $n$ times with doublet $j$ removed,
giving $(K_{(j)}, \tau_{(j)}, L_{(j)})$, and report the standard deviation of each column and of
$L_{(j)}+\tau_{(j)}$. Returns $K$, $\tau$, $L$, $\text{rms} = \sqrt{J/N}$ ($N$ samples), fit \%, $n$, the spreads,
the leave-one-out table, the $J$ map, and the stacked signals for the plots.

\begin{intuition}
If $\tau_{(j)}$ and $L_{(j)}$ jump around while $L_{(j)}+\tau_{(j)}$ stays put, the doublets agree on the
total lateness but each one picks a different point of the valley. The red leave-one-out dots on
\file{cost\_map.png} line up along the valley in that case.
\end{intuition}

\func{video\_delay}{(folder) -> float or None}
Estimates $L_v$ from the yaw doublets. The IMU yaw (telemetry, converted to map convention) does not go through
the video, the ArUco yaw does, so the ArUco yaw lags by $L_v$. With the unwrapped signals $\psi_A$ and $\psi_I$
and, for each yaw window $[a, b]$, a grid $t_i \in [a+0.5, b-0.5)$ of step 0.02\,s:
$$
\hat L_v = \arg\min_{\ell \in \lbrace -0.2, -0.19, \dots, 0.99\rbrace}
\frac{1}{N}\sum_{\text{windows}}\sum_i \Big[\big(\psi_A(t_i+\ell) - \overline{\psi_A}\big) - \big(\psi_I(t_i) - \overline{\psi_I}\big)\Big]^2 ,
$$
the means being taken per window (this removes the unknown heading offset). Requires 50 samples of each and at
least one yaw segment. Then $L_u = L - L_v$.

\func{fit}{(folder, verbose=True) -> dict}
\code{load}, \code{fit\_axis} for each axis, \code{video\_delay}, prints a table
($K \pm \text{sd}$, $\tau$, $L$, $L+\tau$, fit \%, rms, number of doublets), then \code{save}.

\func{save}{(folder, res, Lv, has\_sent)}
Writes \file{model.yaml}, one flow-mapping line per axis:
\begin{lstlisting}
# Tello response per axis, first order + delay: tau dv/dt = K u(t - L) - v   (tello_step_fit)
x: {K: 0.660, tau: ..., L: ..., L_plus_tau: 0.73, sd_K: ..., fit_percent: ..., rms: ..., doublets: 6}
video_delay_L_v: 0.20
\end{lstlisting}
(the values here are illustrative). With matplotlib it also writes \file{step\_fit.png} (per axis: ArUco points,
model, scaled sent command) and \file{cost\_map.png} (per axis: filled contours of $\log_{10}(J/J_{\min})$
clipped at 1, the white $+10\,\%$ contour $J = 1.1 J_{\min}$, the leave-one-out points and the optimum).

\subsection{Validation}

\func{read\_model}{(folder) -> dict}
Parses \file{model.yaml} with PyYAML into \code{\{axis: (K, tau, L)\}}.

\func{latest\_identification}{(before=None) -> Path or None}
The most recent \file{step\_*} folder (the glob \code{step\_*} does not match \code{stepval\_*}) with a
\file{model.yaml}, in the same parent as \code{before} and older than it. The comparison is lexicographic on the
names after replacing \code{stepval} by \code{step}; it works because the timestamp format
\code{\%Y\%m\%d\_\%H\%M\%S} sorts chronologically.

\func{validate}{(ident\_folder, val\_folder, verbose=True) -> str}
For each axis of the validation flight, simulate the identified $(K, \tau, L)$ and the reference
\code{MODEL\_IN\_USE}, fit \emph{only} $c_j, d_j$ (\code{solve(\dots, K=K)}), and report fit \% and RMS for both.
Writes \file{validation.txt} and \file{validation.png}. This is the honest test: the parameters were not tuned on
this data. The Nov 13 numbers quoted in \file{model.py}: 81\,\% forward, 67\,\% left, 87\,\% yaw, against
48/54/30\,\% for the previous model.

\func{main}{()}
Dispatches \code{--validate} or fits the given folder (default: the last \file{step\_2*} in the sessions).

% =====================================================================================================
\section{From the fit to \texorpdfstring{\texttt{model.py}}{model.py}: who uses the numbers}\label{sec:ident-use}

\file{model.py} contains only constants (the drone flies with its protective cage, which makes it slower: $K$
0.66--0.83 instead of $\sim$1.15 without):

\begin{center}
\begin{tabular}{@{}lll@{}}
\toprule
Constant & Value & Meaning \\
\midrule
\code{K\_X}, \code{K\_Y} & 0.66, 0.83 & forward / left gain \\
\code{TAU\_XY} & 0.25 s & horizontal time constant \\
\code{L\_V} & 0.20 s & image taken $\to$ marker pose \\
\code{L\_U} & 0.35 s & command $\to$ motion \\
\code{K\_Z}, \code{TAU\_Z} & 0.9, 0.45 s & vertical (older values, not re-identified) \\
\code{K\_W} & 1.0 & yaw rate gain \\
\code{K\_MEAN} & 0.745 & $(K_X+K_Y)/2$ for direction-free predictions \\
\bottomrule
\end{tabular}
\end{center}

The measured lateness was $L+\tau = 0.73$\,s forward and $0.85$\,s left. Because only that sum is reliable, the
authors fixed one $\tau_{xy} = 0.25$\,s for both axes and chose $L_U$ so that $L_U + L_V + \tau_{xy} = 0.80$\,s
falls in the measured range. Consumers:

\begin{itemize}
\item \textbf{EKF} (\file{ekf\_node.py} passes them to \code{DroneEKF}, \file{ekf\_core.py}). The prediction is
the same model, in the world frame and with the learnt drift $d$:
$\tau\dot{\mathbf v} = R(\psi)\operatorname{diag}(K_x, K_y)\,\mathbf u(t - L_u) + \mathbf d - \mathbf v$,
discretised exactly like \code{response()}: \code{a\_xy = dt / tau\_xy}, $v \mathrel{+}= a(\cdot - v)$. The input
is the buffered \topic{/tello/rc\_sent} read $L_u$ in the past (\code{command\_at}), and each marker pose is applied
at $t_{\text{recv}} - L_v$ and re-propagated to now.
\item \textbf{Position hold} (\file{safety.py}, \code{position\_hold}/\code{hold\_axis}). Three uses:
(i) the error is computed on the position \emph{predicted at the time the command acts},
$e = p^\star - (p + v L_U)$; (ii) when the pilot releases the keys the hold target is where the drone will stop
on its own, $p + (L_U + \tau)v$ (0.60\,s $\times$ speed); (iii) the desired world velocity is rotated to the body
frame and the model is \emph{inverted}, $u_x = b_x / K_X$, $u_y = b_y / K_Y$ (and $/K_Z$ vertically), so a
desired 0.1\,m/s forward becomes a command of 0.15.
\item \textbf{Zone guard} (\file{safety.py}): the position in $H = 1$\,s if the command $w$ is kept, from the
closed-form solution of the model after the delay:
$$
p(H) = p + v L_U + K w (H - L_U) + (v - K w)\,\tau\left(1 - e^{-(H - L_U)/\tau}\right),
$$
with $K = $ \code{K\_MEAN}; obtained by integrating $v(s) = Kw + (v-Kw)e^{-s/\tau}$ over $[0, H-L_U]$.
\item \textbf{Fake drone} (\file{fake\_world.py}): the simulator moves with the same $K_X, K_Y, \tau, L_U$, and
delays the rendered images by $L_V$, so \code{mock:=true} reproduces the real timing.
\item \textbf{Missions} (\file{marker\_tour.py}, \file{patrol.py}, \file{strike.py}): divide their velocity
requests by $K_X$/$K_Y$, and \file{patrol.py} anticipates the heading by $\omega L_U$.
\end{itemize}

\begin{keypoint}
Identification is not a one-off: changing the hardware (e.g.\ flying without the cage) invalidates
\file{model.py}. The procedure is \code{tello\_step\_test}, then \code{schedule:=valid}, then copy the validated
numbers into \file{model.py} (nothing reads \file{model.yaml} automatically).
\end{keypoint}

% =====================================================================================================
\section{Recording every flight: \texorpdfstring{\texttt{session\_recorder.py}}{session\_recorder.py}}

A passive ROS node started by \code{real.launch.py}. It never commands anything; it writes one folder per launch,
\file{YYYYmmdd\_HHMMSS/} in \file{\textasciitilde/tello\_ws/sessions/}:

\begin{center}
\small
\begin{tabular}{@{}lll@{}}
\toprule
File & Content & Source / time base \\
\midrule
\file{frames/<t>.jpg} & annotated image, at most 1 per second, JPEG 85 & \topic{/aruco/debug\_image}, $t$ relative \\
\file{poses.csv} & \code{t, x, y, z, yaw\_deg, sigma\_m} & \topic{/aruco/pose}, $t$ relative \\
\file{telemetry.csv} & telemetry JSON keys + \code{t}, throttled to 2\,Hz & \topic{/tello/telemetry}, relative \\
\file{rc\_sent.csv} & \code{t, cx, cy, cz, cw} sent by the driver & \topic{/tello/rc\_sent}, epoch \\
\file{cmd\_in.csv} & \code{t, cx, cy, cz, cw} asked by the mission/teleop & \topic{/tello/cmd\_vel\_safe}, epoch \\
\file{ekf.csv} & \code{t, x, y, z, sigma\_xy} & \topic{/ekf/pose}, epoch \\
\file{summary.json} & counts, rates, pose range, battery, max temperature & at shutdown \\
\bottomrule
\end{tabular}
\end{center}

``Relative'' means seconds since the recorder started (\code{t0 = time.time()} in the constructor); ``epoch'' is
\code{time.time()} directly. \code{sigma\_m} $= \sqrt{\Sigma_{xx}}$ and \code{sigma\_xy}
$= \sqrt{\max(\Sigma_{xx},\Sigma_{yy})}$. \file{rc\_sent.csv} and \file{ekf.csv} make offline EKF replays possible
(\file{tools/ekf\_replay.py}); \file{cmd\_in.csv} vs \file{rc\_sent.csv} shows what the safety chain changed.

\func{SessionRecorder.\_\_init\_\_}{(self)}
Creates the folder (parameter \code{out\_dir}), opens the CSV files, subscribes (the rc/cmd/EKF writers are
lambdas), and starts a 2\,s timer that flushes every file: when the launch is killed \code{close()} may never run
(``flight 225048 lost its last 40\,s of poses'').

\func{on\_img}{(self, m)}
Counts every debug image, saves one per second: reshapes the raw buffer to $H\times W\times c$, converts RGB to BGR
if needed, writes \file{frames/0012.3.jpg} (name = relative time, format \code{06.1f}).

\func{on\_pose}{(self, m)}
Yaw from the quaternion in degrees; appends the row to memory (for the summary) and to \file{poses.csv}.

\func{on\_tel}{(self, m)}
Keeps one message every 0.5\,s, adds the relative \code{t}, appends to memory and to \file{telemetry.csv}
(header = keys of the first kept message, extra keys ignored).

\func{close}{(self)}
Writes \file{summary.json}: duration, number and rate of debug images, frames saved, number of poses, the ratio
\code{poses / (3 * debug\_images)}, time of the first pose, mean/min/max of $(x,y,z)$, battery at start/end and
maximum \code{temph}; closes the files.

\func{main}{()}
Spins until Ctrl+C or external shutdown, then \code{close()} in a \code{finally}.

Note that \file{step\_test.py} writes its \emph{own} folder (\file{step\_\dots}) with a different
\file{poses.csv} (epoch times, yaw in radians, image stamp): the fit never reads the session recorder's files.

% =====================================================================================================
\section{After the flight: \texorpdfstring{\texttt{flight\_report.py}}{flight\_report.py}}

\code{python3 tools/flight\_report.py [session] [--band 1.7 2.0]} answers ``is the localisation good enough to
fly fast and precisely?'' for one recorder folder (default: the latest \file{2*} folder). It writes
\file{report.txt} and \file{report.png} and ends with a checklist of \code{[OK]}/\code{[CHECK]} verdicts.

\func{read\_csv}{(path) -> list[dict]}
Rows, or \code{[]} for a missing or empty file.

\func{col}{(rows, key, default=nan) -> array}
One column as floats, \code{nan} for missing or unparsable values.

\func{nearest}{(t\_src, v\_src, t) -> (values, dt)}
For sorted source times, the value of the nearest sample at each query time and the time distance $|t_{\text{src}} - t|$;
\code{searchsorted} then a choice between the two neighbours.

\func{folder\_time}{(p) -> float or None}
Epoch time from the last two \code{\_}-separated parts of a folder name (\code{strike\_20261113\_181035} $\to$
the timestamp).

\func{main}{()}
\textbf{Time alignment.} The telemetry rows carry both the recorder's relative \code{t} and the driver's absolute
\code{stamp}; $t_0 = \operatorname{median}(\text{stamp} - t)$ recovers the recorder's start in epoch time, so that
pose times become $t_0 + t$ and can be compared with \file{ekf.csv}. \textbf{Flights}: the rising and falling edges
of \code{flying} give the intervals $[a, b]$; the inner \code{in\_air(t, start=3, end=1)} keeps times inside a flight,
3\,s after takeoff and 1\,s before landing. Then the metrics:

\begin{center}
\small
\begin{tabularx}{\textwidth}{@{}lXl@{}}
\toprule
Metric & Exact computation & Verdict \\
\midrule
Markers in view & 0.05\,s grid over each flight ($a+3$ to $b-1$); a grid point is ``seen'' if a pose exists
within 0.3\,s; share of seen points, longest run of unseen points, poses per second & $\ge 85\,\%$ \\
ArUco vs EKF & for every in-air pose, distance to the nearest EKF sample (if within 0.2\,s): median, 95th
percentile, number of ``jumps'' $> 15$\,cm; EKF $\sigma_{xy}$ median and max & jumps $< 1\,\%$ \\
Height & ArUco $z$ minus the nearest ToF (within 0.6\,s, ToF $> 0.3$\,m); samples within 15\,cm of the median
are ``over the floor'' (furniture reads 0.5--0.9\,m less); median and spread of those & $|\text{med}|, \sigma < 5$\,cm \\
Height band & with \code{--band}: share of EKF $z$ (from 8\,s after takeoff) in the band $\pm 5$\,cm, min/max &
$\ge 95\,\%$ \\
Hover & intervals where the telemetry flag \code{hold} is true, minus the first 2\,s, at least 2\,s and 20 poses
long; distance of each pose to the interval's median point: median, 95th pct; height 95th pct & 95\,\% $< 8$\,cm \\
Pad & takeoff and landing points (\file{set\_home.py}: poses of the first 8\,s of the climb / 2--4\,s before the
landing) vs \file{config/home.yaml} & landing $< 15$\,cm \\
Tracking & every mission folder \file{*\_2*} whose time falls in $[t_0 - 5, \text{last flying}]$ and has a
\file{tracking.csv}: horizontal error $\sqrt{(x_d-x)^2 + (y_d-y)^2}$ RMS and max, height error RMS, per phase &
RMS $< 8$\,cm \\
\bottomrule
\end{tabularx}
\end{center}

The plot shows the EKF position (lines) with the ArUco poses (dots), and the ToF height with the ArUco $z$ and the
band limits, against the time since the first takeoff.

\begin{intuition}
The height check is a free test of the \emph{marker map scale}: if the map were 6\,\% too large (as on Nov 11),
ArUco $z$ minus ToF would not be a constant offset but would vary with the distance to the markers.
\end{intuition}

% =====================================================================================================
\section{Camera calibration: \texorpdfstring{\texttt{calibrate\_camera.py}}{calibrate\_camera.py}}

The localizer's PnP needs the intrinsics. This interactive tool (messages in French) streams the Tello video
(\code{room\_mapper.tello\_stream()}), takes photos automatically of the printed checkerboard (9$\times$6 inner
corners, 25\,mm squares), calibrates with OpenCV and saves \file{config/tello\_camera.yaml}.

\subsection{Camera model}

A point $\mathbf X_c = (X, Y, Z)$ in the camera frame is normalised, $x = X/Z$, $y = Y/Z$, $r^2 = x^2+y^2$,
distorted with the \code{plumb\_bob} (Brown--Conrady) model of 5 coefficients $(k_1, k_2, p_1, p_2, k_3)$:
$$
\begin{aligned}
x_d &= x\,(1 + k_1 r^2 + k_2 r^4 + k_3 r^6) + 2p_1 xy + p_2(r^2 + 2x^2),\\
y_d &= y\,(1 + k_1 r^2 + k_2 r^4 + k_3 r^6) + p_1(r^2 + 2y^2) + 2p_2 xy,
\end{aligned}
$$
then projected with the intrinsic matrix:
$$
\begin{bmatrix} u \\ v \\ 1 \end{bmatrix} = K \begin{bmatrix} x_d \\ y_d \\ 1 \end{bmatrix},\qquad
K = \begin{bmatrix} f_x & 0 & c_x \\ 0 & f_y & c_y \\ 0 & 0 & 1 \end{bmatrix}.
$$
$k_i$ are radial terms (barrel/pincushion), $p_i$ tangential terms (lens not parallel to the sensor).

\subsection{Zhang's method}

\code{cv2.calibrateCamera} follows Zhang (2000). The board is the plane $Z_b = 0$ of its own frame, so for view
$i$ a board point $(X_b, Y_b)$ maps to the (undistorted) image through a homography
$$
s\begin{bmatrix} u \\ v \\ 1\end{bmatrix} = K\begin{bmatrix} \mathbf r_1 & \mathbf r_2 & \mathbf t\end{bmatrix}
\begin{bmatrix} X_b \\ Y_b \\ 1\end{bmatrix} = H_i \begin{bmatrix} X_b \\ Y_b \\ 1\end{bmatrix}.
$$
Since $\mathbf r_1, \mathbf r_2$ are orthonormal, with $H = [\mathbf h_1\ \mathbf h_2\ \mathbf h_3]$ and the symmetric
$B = K^{-\top}K^{-1}$:
$$
\mathbf h_1^\top B\,\mathbf h_2 = 0, \qquad \mathbf h_1^\top B\,\mathbf h_1 = \mathbf h_2^\top B\,\mathbf h_2 .
$$
Two linear equations per view in the 6 entries of $B$ (5 up to scale): three views or more give $B$, hence $K$ by
Cholesky, then each $(R_i, \mathbf t_i)$. This closed form initialises a Levenberg--Marquardt refinement of all
parameters (intrinsics, distortion, poses) minimising the reprojection error
$$
E = \sum_{i}\sum_{j} \big\lVert \mathbf m_{ij} - \pi(K, D, R_i, \mathbf t_i, \mathbf M_j)\big\rVert^2 ,
\qquad \text{RMS} = \sqrt{E / N},
$$
with $N$ = number of views $\times$ 54 corners. The returned \code{rms} is this number in pixels. Varied tilts
are essential: with fronto-parallel views only, the homographies do not constrain the focal length (the
equations become degenerate), which is why the tool insists on tilted views.

\subsection{Functions}

Constants: \code{COLS, ROWS, SQUARE = 9, 6, 0.025}; a $3\times3$ grid of image cells; tilt threshold $20^\circ$;
``still'' $= $ mean corner motion $\le 3$\,px for 0.3\,s; sharpness = variance of the Laplacian on the board
$\ge 30$; 15 to 40 views.

\func{board\_points}{() -> (54, 3) array}
The corners in the board frame: $(0..8, 0..5) \times 0.025$\,m, $Z=0$.

\func{measure}{(pts, size, K0) -> dict}
Describes a detected board: width (fraction of the image width); the board pose by \code{solvePnP} with the
\emph{guessed} intrinsics $K_0$ ($f_0 = (W/2)/\tan(55.6^\circ/2) \approx 910$\,px at 960\,px, the measured FOV) and
no distortion; from the board normal $\mathbf n$ (third column of $R$) the horizontal tilt
$\operatorname{atan2}(n_x, |n_z|)$ and the vertical tilt $\operatorname{atan2}(n_y, |n_z|)$ in degrees; the in-plane
rotation of the first row folded into $[-90^\circ, 90^\circ)$; the set of grid cells touched by corners; and the
normalised centre.

\func{novelty}{(m, cov\_cells, cov\_checks, views) -> (new\_cells, new\_checks, different)}
What a view would add: cells not covered yet, checklist items it satisfies for the first time (near: width
$> 0.45$; far: $< 0.28$; tilted left/right/up/down: beyond $\pm 20^\circ$; diamond: $25^\circ \le |\text{roll}| \le 65^\circ$),
and whether it differs from every photo taken (centre moved $> 0.12$, width $> 0.1$, or a tilt $> 12^\circ$).

\func{draw}{(frame, cov\_cells, cov\_checks, n\_views, status, color, corners, flash, hint="") -> image}
The user interface: green cells, grid lines, detected corners, a white flash on a photo, a top banner, the
checklist (PIL for accented text), the status and hint lines.

\func{verdict}{(rms) -> (word, colour)}
$< 0.35$\,px excellent, $< 0.6$ good, $< 1.0$ medium, otherwise bad.

\func{missing\_hint}{(cov\_cells, cov\_checks) -> str}
Names up to three uncovered cells and the first unchecked item.

\func{run}{(frames\_iter, save=True, show=True, photo\_dir=None) -> (rms, K, D) or None}
Per frame: stop when complete (all 9 cells, $\ge 15$ views, and 5 checks or 25 views) or at 40 views;
\code{findChessboardCorners} (adaptive threshold, normalisation, fast check), \code{cornerSubPix} ($11\times11$
window, 30 iterations or $\epsilon = 0.001$); a photo is taken when the board is sharp, \emph{useful} (new cells or
checks, or merely different while fewer than 15 views) and still for 0.3\,s. C computes (10 views or more), Esc or
closing the window computes if there are 10 views, else quits. Then:
\begin{lstlisting}
rms, K, D, rv, tv = cv2.calibrateCamera(obj, img_pts, size, None, None)
errs = [per-view RMS of cv2.projectPoints(...) - corners]        # pseudo-code
keep = [i for i, e in enumerate(errs) if e < 2.5 * np.median(errs)]
if len(keep) < len(obj) and len(keep) >= 10:
    rms, K, D, _, _ = cv2.calibrateCamera(obj[keep], img_pts[keep], ...)
\end{lstlisting}
i.e.\ one round of outlier rejection on the per-view error $e_i = \sqrt{\tfrac{1}{54}\sum_j \lVert\cdot\rVert^2}$.
A sanity warning is printed if $f_x$ is more than 20\,\% from the expected 910\,px (badly printed board), then
\code{save\_calib} writes the YAML and the last frame is shown undistorted (straight lines must look straight).

\func{main}{()}
\code{--from-photos <dir>} recomputes from saved \file{calib\_*.jpg}; otherwise creates
\file{logs/calibration/<time>/} for the photos and calls \code{run(tello\_stream())}.

\paragraph{Output.} \code{room\_mapper.save\_calib} writes the ROS \code{CameraInfo} YAML layout
(\code{camera\_matrix}, \code{distortion\_model: plumb\_bob}, \code{distortion\_coefficients},
identity \code{rectification\_matrix}, \code{projection\_matrix} $= [K\,|\,0]$) plus the RMS and the date. The
current file (Nov 9, 26 views, $k_3$ fixed to 0 according to its \code{source} field):
$$
\begin{aligned}&f_x = 910.5,\quad f_y = 913.5,\quad c_x = 466.6,\quad c_y = 363.9\ \text{px},\\
&D = (-0.0386,\ 0.0816,\ -0.0013,\ -0.0052,\ 0),\qquad \text{RMS} = 0.24\ \text{px}.\end{aligned}
$$
Check: $2\arctan(480/910.5) = 55.6^\circ$ horizontal FOV, and the distortion is mild.

% =====================================================================================================
\section{Printable markers: \texorpdfstring{\texttt{make\_print\_sheets.py}}{make\_print\_sheets.py}}

Generates PDFs at 300\,dpi ($\text{MM} = 300/25.4 = 11.81$\,px/mm; A4 = $2480\times3508$\,px). The dictionary and
the side come from \file{config/markers.yaml} (\code{DICT\_6X6\_50}, \code{marker\_size: 0.1889}\,m).

\func{font}{(size\_px) -> ImageFont}
First TrueType font found (Arial, DejaVu Sans), else PIL's default.

\func{marker\_page}{(dictionary, marker\_id, side\_mm, dict\_name) -> PIL image}
Draws the marker with \code{cv2.aruco.generateImageMarker} (or \code{drawMarker} for old OpenCV) at exactly
\code{round(side\_mm * MM)} pixels (188.9\,mm $\to$ 2231\,px), centred horizontally. A ``big'' marker (more than
the page width minus 60\,mm) is placed 14\,mm from the top without a cut guide; a smaller one gets a grey cut
rectangle one sixth of the side outside it (the white quiet zone must stay). Below: the ID and ``keep upright'',
the dictionary and size, and a 100\,mm bar to check the print scale with a ruler. Asserts the layout fits; returns a
1-bit image. (The \code{dictionary} argument is unused.)

\func{checkerboard\_page}{(cols=10, rows=7, square\_mm=25) -> PIL image}
A4 landscape, $10\times7$ squares of 25\,mm (295\,px), i.e.\ exactly the $9\times6$ inner corners and 0.025\,m that
\file{calibrate\_camera.py} expects.

\func{main}{()}
Options \code{--markers}, \code{--ids all|lo-hi} (checked against the dictionary size and warned on IDs already in
the map), \code{--paper a4|letter}, \code{--checkerboard}, \code{--out docs/print}. Writes
\file{aruco\_markers\_<ids>\_<paper>.pdf} (one page per marker) and optionally \file{checkerboard\_A4.pdf}.

% =====================================================================================================
\section{Code review notes}

\begin{review}
\begin{enumerate}[leftmargin=*,itemsep=2pt]
\item \file{step\_fit.py:50} \code{MODEL\_IN\_USE} is hard-coded to the Nov 11 model ($K=1.15$, $\tau=0.55$,
$L=0.50$), while \file{model.py:425-430} now has $K_X=0.66$, $K_Y=0.83$, $\tau=0.25$, $L_U+L_V=0.55$. The
validation therefore compares with an old model, not ``the model the EKF / driver use now''. Build it from
\code{tello\_driver.model} instead.
\item \file{step\_fit.py:204-205} The ``sd'' is \code{np.std} of the leave-one-out estimates. The jackknife
standard error is $\sqrt{(n-1)/n\sum_j(\theta_{(j)}-\bar\theta)^2} = \sqrt{n-1}\cdot$\code{np.std}: with 6 doublets
the printed uncertainties are about $2.2\times$ too small.
\item \file{step\_test.py:252} with \file{:452} and \file{:461}: $p_0$ is set to the takeoff hover point, then
\code{recentre()} flies to the zone centre $(0.04, 1.56)$ while \code{safety()} still checks the $\pm0.8$\,m box
around the \emph{takeoff} point. A takeoff more than 0.8\,m from the centre along $x$ or $y$ aborts the test during
that move. Skip the box (or re-anchor it) during the move.
\item \file{step\_test.py:17} The docstring says the drone lands after 3\,s without marker; the default
\code{lost\_timeout} is 15\,s (\file{:101}).
\item \file{step\_test.py:235} The in-flight battery abort threshold (20\,\%) is hard-coded, unlike
\code{min\_battery}.
\item \file{step\_test.py:228} ``ANY KEY = land'' relies on \code{read\_key}, which returns \code{None} for a lone
Esc or an unknown escape sequence (\file{teleop.py:144-148}): Esc does not stop the test.
\item \file{step\_test.py:537} The fit runs only after a clean end; after any \code{Abort} the data are kept but must
be fitted by hand, and \code{tello\_step\_fit} without argument (\file{step\_fit.py:396}) then picks that partial
flight if it is the latest.
\item \file{step\_fit.py:81} The sent command is looked up at the \file{commands.csv} times, i.e.\ the value sent
\emph{before} the driver processed the new request: up to one 50\,ms tick (plus reception jitter) is silently
added to $L$. Building $u$ directly from the \file{sent.csv} timeline would avoid it. Also, the hover phases contain
the closed-loop hold command, so input and noise are correlated there (mild closed-loop bias).
\item \file{step\_fit.py:112} Forward Euler $a=\Delta/\tau$ instead of the exact $1-e^{-\Delta/\tau}$: negligible at
$\tau=0.25$ (2\,\%), 10\,\% at the grid's $\tau=0.05$ (yaw).
\item \file{session\_recorder.py:103} The ratio divides by \code{n\_img * 3}, but the localizer publishes a debug
image every \code{debug\_every\_n = 5} frames by default (\file{aruco\_localizer.py:59}): the summary ratio is
off by $5/3$.
\item \file{session\_recorder.py:61} vs \file{:82} and \file{:92}: \file{ekf.csv}, \file{rc\_sent.csv},
\file{cmd\_in.csv} use epoch time while \file{poses.csv} and \file{telemetry.csv} use time since start, in the same
folder; \file{flight\_report.py:82} must reconstruct $t_0$ from the telemetry \code{stamp}. One time base (epoch)
everywhere would be simpler.
\item \file{session\_recorder.py:91} \code{json.loads} without \code{try} (unlike \file{step\_test.py:174}): one
malformed message kills the recorder's spin. And \file{:95} freezes the CSV header to the first message's keys
(\code{extrasaction="ignore"}): keys added later by the driver are silently dropped.
\item \file{flight\_report.py:127-136} ``Pose jumps against the EKF prediction'' actually compares each pose with the
nearest \emph{published} EKF estimate, which has already fused that pose, at reception time, whereas the pose
describes the image $L_v = 0.2$\,s earlier: at 0.3\,m/s this alone is a 6\,cm bias, and real jumps are partly
absorbed. Comparing at \code{t\_image}, or logging the EKF innovation, would be cleaner.
\item \file{flight\_report.py:204} Missions are searched in the default \code{SESSIONS} folder, not next to the
given session (\code{ses.parent}): with a session path elsewhere the tracking section is silently empty.
\item \file{calibrate\_camera.py:216}, \file{:221}, \file{:263} No \code{cv2.CALIB\_FIX\_K3}, while the current
\file{config/tello\_camera.yaml} was made ``k3 fixed to 0'': a re-run produces a different (and with few wide-field
views, over-fitted) model. The \code{--from-photos} path also skips the outlier rejection.
\item \file{calibrate\_camera.py:219-226} If rejection would leave fewer than 10 views, all views are kept but the
message still prints \code{len(keep)} as the number used.
\item \file{make\_print\_sheets.py:4,8} The docstring says the default run writes the checkerboard and that the
marker is 20\,cm; the checkerboard needs \code{--checkerboard} (\file{:87}) and the size comes from
\file{markers.yaml} (0.1889\,m). The \code{dictionary} parameter of \code{marker\_page} (\file{:34}) is unused.
\end{enumerate}
\end{review}
